\section{Related Work}
\label{sec:related}

We position our work within four lines of research:
verifiable ML inference, dedicated matrix multiplication verification,
code-based proof systems, and the sumcheck approach.

\paragraph{Verifiable ML inference.}
SafetyNets~\cite{ghodsi2017safetynets} was among the first to propose
interactive proofs for deep neural networks, employing the sumcheck protocol
to verify matrix multiplications layer by layer.
zkCNN~\cite{liu2021zkcnn} extended this to convolutional layers by reducing
convolution to polynomial multiplication and leveraging the
GKR protocol~\cite{goldwasser2015delegating}.
vCNN~\cite{lee2024vcnn} further optimized the verification of quantized
convolution operations.
Mystique~\cite{weng2021mystique} and zkML~\cite{chen2024zkml} have explored
practical zero-knowledge proof systems for general ML inference, including
Transformer architectures.
In all of these systems, matrix multiplication is the dominant cost center.
Our protocol is complementary: it provides a drop-in replacement for the
matrix multiplication verification subroutine that reduces the SNARK circuit
size from $\mathcal{O}(k^2)$ to $\mathcal{O}(tk)$.
Composing $\protocol$ across the layers of a neural network is a natural
extension that we outline in Section~\ref{sec:conclusion}.

\paragraph{Dedicated matrix multiplication verification.}
zkMatrix~\cite{cong2024zkmatrix} proposed a batched short-proof protocol for
committed matrix multiplication using customized polynomial commitments,
achieving proof sizes of $\mathcal{O}(\sqrt{k})$ group elements.
DualMatrix~\cite{cong2026dualmatrix} combined polynomial evaluation and
inner-product arguments to reduce the prover cost, though the prover still
performs $\mathcal{O}(k^2)$ work inside the proof system.
zkVC~\cite{zhang2025zkvc} introduced Constraint-Reduced Polynomial Circuits
(CRPC) and Prefix-Sum Query (PSQ) techniques to reduce the R1CS constraint
count, reporting a $12\times$ prover speedup and demonstrating applicability
to Transformer inference.
Bennett et al.~\cite{bennett2023matrix} explored the use of coding theory for
matrix multiplication verification in an information-theoretic setting,
establishing theoretical connections between code distance and verification
complexity.
Our work shares the coding-theoretic perspective
of~\cite{bennett2023matrix} but differs in two key respects:
(i)~we operate in the computational setting with Merkle commitments and
SNARKs, yielding a \emph{non-interactive} protocol with succinct proofs;
and (ii)~we provide a concrete implementation with experimental evaluation,
demonstrating practical gains over the Freivalds baseline.
Relative to zkMatrix, DualMatrix, and zkVC, our approach avoids
polynomial commitment machinery entirely, relying instead on
hash-based commitments that require no trusted setup.

\paragraph{Code-based proof systems.}
Ligero~\cite{ames2017ligero} pioneered the use of interleaved Reed--Solomon
codes to construct transparent zero-knowledge proofs with
$\mathcal{O}(\sqrt{n})$ proof size.
Brakedown~\cite{golovnev2023brakedown} introduced expander-based linear codes
as a drop-in replacement for polynomial commitments, achieving linear-time
prover complexity for general-purpose R1CS.
Orion~\cite{xie2022orion} combined code-based commitments with the sumcheck
protocol.
FRI~\cite{ben2018fast} (Fast Reed--Solomon Interactive Oracle Proofs of
Proximity) established the foundation for proximity testing of
Reed--Solomon codes, underpinning the STARK family of proof systems.
Our work draws on the code-based commitment philosophy of Brakedown but
applies it to the \emph{structured} setting of matrix multiplication, where
the linearity of encoding interacts directly with the Freivalds reduction
to yield a circuit-size improvement beyond what a generic code-based SNARK
achieves.
Specifically, while Brakedown provides a general $\mathcal{O}(n)$-prover
SNARK for any circuit of size $n$, na\"ively feeding the
$\mathcal{O}(k^2)$ Freivalds circuit into Brakedown still yields an
$\mathcal{O}(k^2)$ prover.
$\protocol$'s contribution is to \emph{restructure the problem itself}---via
the relation chain
$\mathcal{R}_{\mathsf{orig}} \to \mathcal{R}_{\mathsf{Freivalds}} \to \mathcal{R}_{\mathsf{SNARK}}$---so
that the circuit presented to the SNARK is only $\mathcal{O}(tk)$.

\paragraph{Sumcheck-based approaches.}
The sumcheck protocol~\cite{goldwasser2015delegating} is a powerful
interactive proof technique for verifying the evaluation of multivariate
polynomials.
When applied to matrix multiplication, it reduces the verifier's work to
evaluating multilinear extensions of the input matrices at a random point,
which requires $\mathcal{O}(k^2)$ field operations for the verifier (or
$\mathcal{O}(k^2)$ oracle queries).
SafetyNets~\cite{ghodsi2017safetynets} and
zkCNN~\cite{liu2021zkcnn} use this approach.
Unlike sumcheck-based methods, our proximity-based checks are purely
combinatorial: they compare codeword coordinates via inner products rather
than evaluating polynomial extensions, avoiding the algebraic overhead
of sumcheck rounds.
Moreover, our protocol naturally interfaces with any SNARK backend (the
local checks are expressed as R1CS or PlonK constraints), whereas
sumcheck-based systems typically require a specialized verifier or a
dedicated GKR-to-SNARK compilation step.
